% ============================================================
% Chapter 2: Literature Review
% ============================================================

\chapter{Literature Review}
\label{chap:litrev}

This chapter surveys the body of work that contextualises, motivates, and
directly informs the present study. The review is organised into six thematic
sections: the emergence and definition of AI coding agents
(Section~\ref{sec:lr:agents}); empirical studies of adoption rates and trends
on GitHub (Section~\ref{sec:lr:adoption}); detection methodology and the
multi-signal challenge (Section~\ref{sec:lr:detection}); how developers
interact with and adopt AI tools (Section~\ref{sec:lr:interaction}); the
JavaScript, TypeScript, and npm ecosystem as a research subject
(Section~\ref{sec:lr:ecosystem}); and geographic dimensions of open-source
software development (Section~\ref{sec:lr:geography}). The chapter concludes
by synthesising the identified research gaps that this thesis addresses
(Section~\ref{sec:lr:gaps}).

% ------------------------------------------------------------
\section{AI Coding Agents: Emergence and Definitions}
\label{sec:lr:agents}
% ------------------------------------------------------------

The term \emph{coding agent} refers to an AI system embedded in an autonomous
loop that plans a sub-task, executes an action (writing or modifying code,
running tests, reading documentation), observes the result, and iterates
until a goal is reached~\cite{robbesAgenticMuch2026}. This distinguishes
agents from earlier AI coding \emph{assistants}, which provide
suggestions inline but require the human to accept each change individually.
The distinction is operationally significant: an assistant leaves a trace
only when the developer accepts a suggestion, whereas an agent may author
entire functions, write tests, and open pull requests autonomously, leaving
a richer and more diverse set of observable artefacts.

The practical frontier between assistants and agents is blurred by how tools
are configured. GitHub Copilot, for example, operates as an autocomplete
assistant in its default IDE mode but transitions to an agentic mode when
given access to a terminal or a repository via Copilot Workspace. Similarly,
Cursor and Windsurf expose both completion and multi-step agentic modes
within the same interface. Chatlatanagulchai et al.~\cite{chatlatanagulchai2025}
document the growth of \emph{context files}---configuration artefacts such as
\texttt{CLAUDE.md}, \texttt{.cursorrules}, and \texttt{AGENTS.md}---which
signal that a developer has configured an agent to operate with project-level
instructions rather than just accepting inline suggestions. The presence of
such files is the most reliable publicly observable signal that agentic (as
opposed to merely AI-assisted) development is taking place.

The key tools covered by this study and their approximate launch dates are
listed in Table~\ref{tab:tools}. The most consequential for the npm ecosystem
is Claude Code (launched March 2025), which targets command-line-first
developers and stamps every commit it authors with a
\texttt{Co-Authored-By: Claude <noreply@anthropic.com>} trailer by
default~\cite{robbesHeuristics2026}. This default behaviour makes Claude Code
the most detectable agent in repository mining studies.

\begin{table}[h]
\centering
\caption{Key AI coding tools and their approximate launch dates.}
\label{tab:tools}
\begin{tabular}{llll}
\hline
\textbf{Tool} & \textbf{Vendor} & \textbf{Mode} & \textbf{Launch} \\
\hline
GitHub Copilot       & GitHub/Microsoft & Assistant + Agent & Jun 2021 \\
Cursor               & Anysphere        & Assistant + Agent & 2023 \\
Aider                & Paul Gauthier    & Agent (CLI)       & 2023 \\
Devin                & Cognition AI     & Agent (full)      & Mar 2024 \\
Claude Code          & Anthropic        & Agent (CLI)       & Mar 2025 \\
Windsurf             & Codeium          & Assistant + Agent & Nov 2024 \\
Gemini Code Assist   & Google           & Assistant + Agent & 2024 \\
\hline
\end{tabular}
\end{table}

Galster et al.~\cite{galsterConfig2026} conducted a large-scale empirical
analysis of configuration artefacts across 4,738 repositories, finding
148,519 AI-co-authored commits and 15,591 configuration artefact files.
TypeScript was the most represented primary language (24.5\% of repos in their
sample), providing early evidence that the npm ecosystem is a particularly
active site of AI agent adoption. A companion study by the same group
\cite{galsterHarness2026} characterised the structure of agent harness files,
finding that 85.5\% of skill definitions lack dynamic scripts, revealing that
much agentic configuration is declarative rather than programmatic.

% ------------------------------------------------------------
\section{Empirical Studies of Coding Agent Adoption}
\label{sec:lr:adoption}
% ------------------------------------------------------------

The two most comprehensive empirical studies of coding agent adoption on
GitHub are the \emph{Agentic Much?}~\cite{robbesAgenticMuch2026} and
\emph{Agentic Very Much!}~\cite{robbesAgenticVery2026} papers by Robbes
et al., which together establish the baseline adoption rates against which
this thesis measures the JavaScript/TypeScript ecosystem specifically.

\subsection{Adoption in Established Repositories}

Robbes et al.~\cite{robbesAgenticMuch2026} studied 128,018 GitHub repositories
across all languages using a snapshot dated 21 February 2026. Applying a
taxonomy of 196 heuristics covering 63 agents, they found an overall adoption
rate of \textbf{22.20\%--28.66\%}, depending on which heuristic set is used.
The 6.46 percentage-point gap between these bounds represents repositories
that are missed when a narrower signal set is applied---the \emph{undercounting
delta} that this thesis explicitly quantifies for the JS/TS ecosystem.

Beyond the headline rate, Robbes et al. identified several patterns that
motivate the design of the present study. First, adoption is positively
correlated with repository size, activity, and star count, which justifies
the size-stratified sampling strategy adopted in Stage~1 of this thesis.
Second, AI-generated commits are \textbf{3$\times$ larger} than human commits
(median 31 vs.\ 11 added lines), suggesting that agents introduce code in
larger blocks than individual autocomplete acceptances. Third, the most common
task type for AI-authored commits is feature work (35.7\%), followed by bug
fixing (29.9\%). Fourth, temporal analysis reveals a rapid acceleration of
adoption from mid-2025 onwards, coinciding with the release of Claude Code
and the wider availability of agentic modes in Cursor and Copilot.

A crucial limitation of this study, acknowledged by the authors, is that the
cross-language aggregation conceals ecosystem-specific patterns. The 22--28\%
figure pools together low-adoption ecosystems (e.g.\ legacy C repositories
maintained by researchers) with potentially high-adoption ecosystems such as
npm. This aggregation gap is the primary scope motivation for the present
thesis.

\subsection{Adoption in New Repositories}

A companion study by the same group~\cite{robbesAgenticVery2026} restricted
the analysis to 12,794 repositories \emph{created} after 29 August 2025---the
point at which the researchers consider the agentic AI era to have matured.
This design choice directly parallels the corpus construction approach of the
present thesis, which requires \texttt{created:>=2025-02-24}.

The results are striking: adoption in new repositories reaches
\textbf{71.83\%--76.15\%}, roughly three times the rate in older
repositories. The median commit agent ratio (fraction of commits authored or
co-authored by an agent) is approximately 30\% in new repositories versus
10\% in older ones. The proportion of repositories in the ``pervasive
adoption'' category---where agents contribute to more than 50\% of
commits---rises from 21.1\% in older repositories to 41.2\% in new ones.
Claude Code accounts for approximately three times more signals than Copilot
in new repositories, consistent with its more aggressive default attribution
behaviour and its CLI-first design that integrates naturally with terminal
workflows.

These findings motivate the creation date filter applied in Stage~1 of this
thesis. Repositories created before the agentic era were built without agents
from day one; including them in a study of agent adoption conflates
non-adoption with \emph{pre-adoption}, inflating the apparent non-adoption
rate. By focusing on repositories created on or after 24 February 2025,
the present study measures adoption among projects that had access to
mature agentic tools throughout their entire development lifecycle.

\subsection{A Census-Scale Study}

Khosravani and Mockus~\cite{khosravani2026} provide the largest-scale
analysis to date, surveying over 180 million GitHub repositories using
four signal types: bot account authorship (Type~A), commit message trailers
(Type~B), author email patterns (Type~C), and configuration file presence
(Type~D). Their central finding is that single-signal bot-account detection
misses \textbf{up to 30$\times$} the volume of Claude Code commits detected
by the full multi-signal approach---the strongest empirical evidence for why
any single-signal detection methodology produces a severe lower bound.
Their study also documents 320,000+ AI-authored commits per month as of early
2026, confirming that agent adoption has crossed from an experimental to an
industrial phenomenon.

% ------------------------------------------------------------
\section{Detection Methodology and the Multi-Signal Challenge}
\label{sec:lr:detection}
% ------------------------------------------------------------

Translating the conceptual question ``does this repository use AI coding
agents?'' into an operational answer requires observable signals. The
literature converges on four families of signals~\cite{robbesHeuristics2026,
galsterConfig2026, khosravani2026}:

\begin{enumerate}
    \item \textbf{Configuration artefacts:} Files at the repository root that
    configure an agent's behaviour---\texttt{CLAUDE.md}, \texttt{AGENTS.md},
    \texttt{.cursorrules}, \texttt{.aider.conf.yml}, \texttt{.roomodes}, and
    others. Galster et al.~\cite{galsterConfig2026} catalogue 15,591 such files
    across their corpus. These signals have high precision (a \texttt{CLAUDE.md}
    file is nearly unambiguous evidence of intentional Claude Code use) but
    limited recall (many teams use agents without creating configuration files).

    \item \textbf{Commit message trailers:} The \texttt{Co-Authored-By:}
    trailer is injected by several agents by default, most notably Claude Code
    (\texttt{noreply@anthropic.com}). Robbes et al.~\cite{robbesHeuristics2026}
    document 2.7 million commits carrying this email address and catalogue the
    full set of trailer patterns used by 63 agents. Trailer-based detection
    has high precision but is easily suppressed by developers who clean their
    commit history.

    \item \textbf{Pull request markers:} Agents that operate in PR-creation
    mode (Copilot, Devin, Codex) leave markers in PR bodies, branch names
    (\texttt{copilot/fix-*}, \texttt{codex/*}), and PR labels. Robbes
    et al.~\cite{robbesHeuristics2026} document 2.1 million branches with the
    \texttt{codex/} prefix. These signals are PR-scope only and miss
    commit-level agent usage.

    \item \textbf{Bot account authorship:} Some agents operate as registered
    GitHub apps (e.g.\ \texttt{github-copilot[bot]},
    \texttt{devin-ai-integration[bot]}), making their commits attributable
    directly via the author login. This is the highest-precision signal but
    covers only a minority of agent types.
\end{enumerate}

Ghaleb~\cite{ghaleb2026} approaches detection from a machine-learning
perspective, training an XGBoost/Random Forest classifier on 33,580 pull
requests with 97.2\% F1-score. The classifier's key features include
multiline commit density (Codex: 67.5\%), PR description length (Copilot:
38.4\%), bullet-point formatting (Cursor: 17.2\%), and conditional statement
density (Claude Code: 27.2\%). While achieving high accuracy, this approach
requires labelled training data and is therefore complementary rather than
a direct substitute for heuristic-based detection in a new ecosystem.

The relationship between signal types has direct methodological implications.
Mao et al.~\cite{mao2026} demonstrate that restricting detection to a single
signal type---the approach of most pre-2025 studies---systematically
undercounts adoption. The present thesis operationalises this insight by
reporting both a \emph{single-signal} rate (commit trailers only, for
comparability with prior work) and a \emph{multi-signal} rate (all signal
types), treating the gap between them as an estimate of the undercounting
delta.

The study also distinguishes between two adoption levels, following the
conceptual framework implicit in Galster et al.~\cite{galsterConfig2026}
and Chatlatanagulchai et al.~\cite{chatlatanagulchai2025}. Repositories with
sparse, isolated signals (a single commit trailer, no configuration file) are
classified as \emph{AI-assisted}---the developer may have used the tool
occasionally without integrating it into their workflow. Repositories where
a configuration file is present, or where multiple signal types co-occur, or
where signals appear with high frequency, are classified as
\emph{agentic}---indicating systematic integration of the agent into the
development process.

% ------------------------------------------------------------
\section{Developer Interaction with AI Coding Tools}
\label{sec:lr:interaction}
% ------------------------------------------------------------

Understanding why some repositories produce observable agent signals while
others do not requires insight into how developers interact with these tools
in practice.

\subsection{Bimodal Interaction Patterns}

Barke et al.~\cite{barke2023} conducted a grounded theory study of 20
programmers using GitHub Copilot on structured coding tasks, producing the
first theoretically grounded account of developer--AI interaction. They
identified a stable bimodal pattern: \emph{acceleration mode}, in which
developers know what to write and use Copilot to get there faster (code
prompts dominant: 71.4\%), and \emph{exploration mode}, in which developers
are uncertain and use Copilot to discover API usage and problem structure
(comment prompts dominant: 57.2\%). Exploration interactions take more than
twice as long as acceleration interactions (248.6 vs.\ 104.7 minutes total
across participants). The dominant determinant of which mode a developer
enters is task familiarity: familiar logic drives acceleration, unfamiliar
APIs drive exploration.

This bimodal model has two implications for the present study. First,
exploration-mode interactions---where developers write natural-language
comments as prompts and accept large generated blocks---are more likely to
leave detectable traces in commit diffs and PR descriptions. Second, the
model explains why adoption rates may differ across developer communities:
developers who work predominantly in unfamiliar territory (e.g.\ students
or developers in a new ecosystem) may produce richer agent traces than
experienced developers in acceleration mode, complicating the interpretation
of raw signal counts.

\subsection{Productivity and Acceptance Rates}

Ziegler et al.~\cite{ziegler2022} conducted the largest empirical study of
AI coding tool productivity to date, matching survey responses from 2,631
GitHub Copilot users to four weeks of IDE telemetry. The overall acceptance
rate was \textbf{27\%}, with users accepting more than 31 suggestions per
active day on average. Acceptance rate (\texttt{accepted\_per\_shown}) was
the strongest single predictor of self-reported productivity ($r = 0.24$,
$p < 0.0001$), though the correlation left substantial unexplained variance,
underscoring that productivity is a multidimensional construct.

Two language-specific findings are directly relevant to this thesis.
JavaScript achieved a higher acceptance rate (30.8\%) than TypeScript
(26.3\%)---consistent with JavaScript's looser type system reducing the
frequency of suggestion rejection due to type mismatches. This difference
implies that a JS/TS-focused study should expect somewhat more AI-generated
code per active user in JavaScript repositories than TypeScript repositories,
all else equal. Acceptance rates also varied with coding context: during
typical working hours they were 21.2\%, rising to 23.5\% on weekends---a
reminder that repository-level signals integrate across heterogeneous
developer time regimes.

Subsequent work by Mazloomzadeh et al.~\cite{mazloomzadeh2026} and
Ogenrwot and Businge~\cite{ogenrwot2026} extended the productivity and quality
evidence to the pull-request level, finding that AI-assisted PRs are accepted
at rates comparable to human PRs (83.8\% vs.\ 91.0\%) but that AI-authored
commits are reverted more frequently (Copilot: 7.6\%, Claude Code: 7.2\%) than
human-authored commits (4.7\%)~\cite{oukhay2026}. Watanabe et al.\
\cite{watanabe2025} found that AI assistance increases the volume of code
review comments without proportionally increasing review quality, suggesting
that reviewers apply additional scrutiny to suspected AI-generated code.

% ------------------------------------------------------------
\section{The JavaScript, TypeScript, and npm Ecosystem}
\label{sec:lr:ecosystem}
% ------------------------------------------------------------

The selection of the JavaScript/TypeScript/npm ecosystem as the unit of
analysis is motivated by three converging lines of evidence: ecosystem scale,
tool affinity, and research gap.

\subsection{Ecosystem Scale and Growth}

The GitHub Octoverse 2025 report~\cite{octoverse2025} documents that
TypeScript overtook Python as the most widely used programming language on
GitHub in August 2025, with 2,636,006 monthly contributors---a year-on-year
increase of 66.6\%. The report explicitly attributes this growth to AI coding
agent adoption, noting that TypeScript's strict typing and comprehensive IDE
tooling make it particularly well-suited for agentic code generation, which
benefits from type information to avoid runtime errors. JavaScript remains
the largest language by raw repository count; together, JavaScript and
TypeScript account for the plurality of active GitHub repositories.

The npm registry, which serves as the package management infrastructure for
both languages, is the world's largest software package registry.
Chatzidimitriou et al.~\cite{chatzidimitriou2018} characterised the npm
ecosystem through the npm-miner infrastructure, analysing 598,862 packages
and finding that quality metrics---including cyclomatic complexity, code
duplication, and security vulnerabilities---follow heavily skewed power-law
distributions, with extreme outliers driven by the registry's openness. The
mean npm package has 4.01 direct dependencies and 1,577 lines of code,
providing a quantitative baseline for what a ``typical'' npm repository looks
like. These characteristics motivate the lockfile filter applied in this
thesis: a repository containing \texttt{package.json} alongside a lockfile
(\texttt{package-lock.json}, \texttt{yarn.lock}, or \texttt{pnpm-lock.yaml})
demonstrates active npm dependency management, distinguishing it from
repositories where npm is incidental to the project's primary language.

\subsection{Ecosystem Affinity with AI Coding Agents}

Evidence from configuration artefact studies confirms that the npm ecosystem
has disproportionate affinity with AI coding agents. Galster et al.\
\cite{galsterConfig2026} found TypeScript to be the most common primary
language (24.5\%) in repositories containing agent configuration files, while
a companion harness engineering study~\cite{galsterHarness2026} found
TypeScript to be the top language (27.8\%) in repositories with agent skill
definitions. Robbes et al.~\cite{robbesAgenticVery2026} found Claude Code to
be dominant in new-project repositories by a factor of 3$\times$ over
Copilot---Claude Code's CLI-first design integrates naturally with the npm
toolchain, which is itself terminal-centric.

Markov\'{a}~\cite{markova2025} documents a migration trend from JavaScript
to TypeScript in existing repositories, driven partly by the desire to
improve AI coding tool performance: typed codebases produce fewer incorrect
suggestions and require less correction. This migration pattern means that
the JS/TS corpus is not static---some repositories in the study window will
be mid-migration, which the \texttt{primary\_language} column in the corpus
captures.

\subsection{GitHub as a Data Source}

Constructing a JS/TS/npm corpus from GitHub requires careful attention to
the data quality issues documented by Kalliamvakou et al.\
\cite{kalliamvakou2016} in their landmark study of the promises and perils of
mining GitHub. Analysing 6.8 million public repositories, they found that
67\% have only a single committer, the median project lifespan is 9.9 days,
and only 12.5\% were active in the preceding month. Furthermore, 63.4\% of
projects in a random sample were genuine software development efforts, while
the remainder were experimental, storage, or academic repositories. Only
10\% of collaborative repositories ever used pull requests, and approximately
40\% of pull requests that are truly merged appear as non-merged in the GitHub
API due to direct-push merges.

These findings directly inform the corpus construction strategy of the present
thesis. The requirements of at least four stars (visibility filter), at least
50 commits (activity filter), creation after 24 February 2025 (temporal
filter), and non-fork status (originality filter) together operationalise the
recommendations of Kalliamvakou et al.\ for constructing a corpus of
active, original software development repositories. Dabić et al.\
\cite{dabic2021} provide additional methodological grounding: the SEART GHS
sampling tool, which underpins several prior MSR studies, applies similar
filters and confirms that star-based stratification produces corpora with
stable cross-study comparability properties.

% ------------------------------------------------------------
\section{Geographic Dimensions of Open-Source Software Development}
\label{sec:lr:geography}
% ------------------------------------------------------------

The geographic comparison between Egyptian and global JS/TS developers
constitutes the second research question of this thesis. The relevant
literature spans geography of contribution, geographic attribution
methodology, and the absence of prior community-specific adoption studies.

\subsection{Geographic Distribution of OSS Contribution}

Wachs et al.~\cite{wachs2022} conducted the most comprehensive empirical
study of geographic patterns in open-source software on GitHub, analysing
contributor location data for millions of repositories. Their central finding
is that OSS contribution is strongly concentrated in North America and Western
Europe, with developers in these regions contributing to disproportionately
high-visibility repositories. They document a measurable correlation between
a country's GDP per capita and the visibility of its developers' OSS
contributions.

Crucially for the methodology of this thesis, Wachs et al.\ validate GitHub
profile location self-reporting as a practical (if imperfect) basis for
geographic attribution. They find that self-reported location is broadly
consistent with other geographic signals (employer IP addresses, conference
attendance records) for the population of developers who fill in their
location field, though they caution that the blank-location population---who
cannot be attributed---introduces systematic undercoverage of
lower-income-country developers. This limitation directly motivates the
snowball expansion component of Stage~2 of this thesis, which recovers
Egyptian developers who are socially connected to location-tagged Egyptian
accounts but do not themselves report a location.

\subsection{The Egyptian Developer Community}

Egypt is the most populous Arabic-speaking country and has a rapidly growing
technology sector, with several universities (Cairo University, the American
University in Cairo, and the German University in Cairo, among others)
producing large annual cohorts of software engineering graduates. Egyptian
developers represent a case of particular interest for AI adoption studies
because they combine:
(a) a predominantly English-language professional coding environment
(reducing language-barrier confounds that would complicate studies of, e.g.,
Chinese-language developers);
(b) an active GitHub presence documented by the Stage~2a location search,
which identified 1,876 seed accounts through explicit self-reported
Egyptian location strings in both English (\textit{Egypt}, \textit{Cairo},
\textit{Alexandria}, \textit{Giza}) and Arabic (\textit{مصر},
\textit{القاهرة}, \textit{الإسكندرية}, \textit{الجيزة}); and
(c) a development context characterised predominantly by personal, academic,
and small-scale commercial projects rather than high-visibility open-source
libraries, as evidenced by the strong skew of the Egyptian sub-corpus toward
low-star repositories (OBS-001).

No prior study has examined AI coding agent adoption specifically within the
Egyptian or MENA developer community. The closest work is Wachs
et al.~\cite{wachs2022}, which documents the geographic distribution of OSS
contribution but does not address tool adoption. This thesis is therefore the
first to test whether the 22--76\% adoption range documented at the global
level holds within a specific national developer community, and to investigate
the mechanisms---including team composition and international collaboration
exposure---that might explain community-level differences.

% ------------------------------------------------------------
\section{Research Gaps}
\label{sec:lr:gaps}
% ------------------------------------------------------------

The foregoing review identifies three research gaps that the present thesis
directly addresses.

\textbf{Gap 1 --- The Ecosystem Scope Gap.}
All major empirical studies of coding agent adoption (Robbes et al.\
\cite{robbesAgenticMuch2026,robbesAgenticVery2026}, Khosravani and Mockus
\cite{khosravani2026}) report cross-language aggregate rates. No study has
measured adoption specifically within the JavaScript/TypeScript/npm ecosystem,
despite this ecosystem being the most active site of agent adoption according
to configuration artefact evidence~\cite{galsterConfig2026,octoverse2025}.
The aggregate rates therefore provide a lower bound for what should be
expected within npm, but cannot answer whether JS/TS developers adopt
more, less, or differently than the general GitHub population.

\textbf{Gap 2 --- The Undercounting Gap.}
The 6.46 percentage-point gap between the narrow and broad heuristic sets in
Robbes et al.~\cite{robbesAgenticMuch2026}, and the 30$\times$ undercounting
factor documented by Khosravani and Mockus~\cite{khosravani2026}, establish
that single-signal detection severely misrepresents true adoption. However,
no study has quantified this undercounting delta specifically for the
JavaScript/TypeScript ecosystem, where the package manager infrastructure
(npm dependencies as signals) provides a fifth signal type not applicable to
other ecosystems.

\textbf{Gap 3 --- The Community Comparison Gap.}
Prior adoption studies report blended figures that obscure community-level
variation. The geography of OSS contribution is well-documented
\cite{wachs2022}, but the intersection of geographic community membership
and AI coding agent adoption has not been studied. In particular, no prior
work has investigated whether the adoption gap between high-income and
lower-income developer communities---predicted by the Wachs et al.\ framework
and motivated by the cost and discoverability of agentic tools---is observable
in repository-level adoption signals.

Together, these gaps frame the contributions of the present thesis, which
quantifies adoption within the JS/TS/npm ecosystem (addressing Gap~1), reports
both single-signal and multi-signal rates to estimate the undercounting delta
(Gap~2), and compares Egyptian and global adoption rates stratified by
repository size (Gap~3).
